\documentclass[10pt,a4paper]{amsart}
\usepackage[margin=19mm]{geometry}
\usepackage{amsmath,amssymb,mathtools}
\usepackage[scr=rsfs,cal=euler]{mathalfa}
\usepackage{xfrac}
\usepackage[unicode,hidelinks,bookmarks=false]{hyperref}
\newcommand{\ie}{i.e.\ }
\newcommand{\cf}{cf.\ }
\newcommand{\resp}{respectively\ }
\newcommand{\Subsection}{\subsection}
\providecommand{\nobreakdash}{\nobreak\hbox{-}}
\setlength{\emergencystretch}{2em}
\multlinegap0pt
\usepackage{amssymb}
\usepackage{mathtools}
\usepackage[shortlabels]{enumitem}
\usepackage{tikz}
\newtheorem{theorem}{Theorem}
\newtheorem{lemma}{Lemma}
\newcommand{\bbbn}{\mathbb N}
\newcommand{\bbbr}{\mathbb R}
\newcommand{\eps}{\varepsilon}
\title{Differentiability of Lipschitz mappings into metric spaces: area and co-area formulas}
\author{Behnam Esmayli, Pawe\l{} Goldstein, Piotr Haj\l{}asz}
\date{Source: arXiv:2609.09121v1 (CC BY 4.0)}
\begin{document}
\maketitle

\noindent\emph{Source and licence.} Differentiability of Lipschitz mappings into metric spaces: area and co-area formulas. Version arXiv:2609.09121v1 (CC BY 4.0), distributed by arXiv under the Creative Commons Attribution 4.0 International licence, http://creativecommons.org/licenses/by/4.0/, as recorded in the arXiv licence metadata for this exact version and shown on the abstract page https://arxiv.org/abs/2609.09121v1. Full licence notice: \texttt{licenses/arxiv-2609.09121-CC-BY-4.0.txt}.\par\smallskip

\noindent\textit{Editorial scope.} Counted unit: the article's theorem BT5 (source0.tex lines 694--702). The article's complete original proof of it is delivered with it (source0.tex lines 724--821, one \texttt{proof} environment of 98 lines, which the article places away from the statement). No mathematics is written, repaired or re-derived: the statement and the proof are the article's own lines, in the article's own notation, and no retained line is substituted or re-flowed.  Retained uncounted as the source-local context the proof needs: lines 704--721.  Nothing was omitted from any retained span, so the package carries no omission record.  No retained line contains a citation, so the wrapper carries no bibliography and no bibliography entry.  No retained line contains a figure environment, an image-inclusion command or drawing code, so no image is delivered and the package declares no asset dependency.  Editorial formatting only, in addition: an AMS \texttt{amsart} preamble that restates the article's own declarations and macros used by the retained lines, namely the environments \texttt{theorem}, \texttt{lemma} on separate counters, together with the standard packages those lines need.  The retained environments print in this document as Theorem 1, Lemma 1 in order of appearance, whereas the article prints the same items with the numbering of its own full document.  The complete native source of the article as distributed in the arXiv e-print for this exact version is bundled unchanged as \texttt{sources/209780/source0.tex}.
\begin{theorem}[Whitney]
\label{BT6}
Let $K\subset\bbbr^n$ be a compact set and let $f:K\to\bbbr$, $L:K\to\bbbr^n$
be continuous functions. Then there is a function $F\in C^1(\bbbr^n)$ such that
$$
F|_K=f
\qquad
\mbox{and}
\qquad
DF|_K=L
$$
if and only if
$$
\lim_{{\scriptstyle x,y\in K,\, x\neq y}\atop {\scriptstyle |x-y|\to 0}}
\frac{|f(y)-f(x)-L(x)(y-x)|}{|y-x|} = 0.
$$

\end{theorem}

\begin{theorem}[Federer]
\label{BT5}
If $f:\Omega\to\bbbr$ is a Lipschitz continuous function defined on an open set $\Omega\subset\bbbr^n$, then for every $\eps>0$ there is an $f_\eps\in C^1(\bbbr^n)$
and a closed set $E\subset\Omega$ such that $|\Omega\setminus E|<\eps$ and
$$
f(x)=f_\eps(x)\text{ and } Df(x)= Df_\eps(x) \text{ for all } x\in E.
$$
Moreover, if $|f|\leq M$ in $\Omega$, then we can take $f_\eps$ with $|f_\eps|\leq 2M$ in $\bbbr^n$.
\end{theorem}

\begin{proof}[Proof of Theorem~\ref{BT5}]
Let $L(x)=Df(x)$. By Rademacher's theorem, $L$ is defined a.e.\ and it is a measurable function.

Assume for a moment that $f:\bbbr^n\to\bbbr$ is defined on $\bbbr^n$ and that $f$ has compact support in an open ball $B$. According to the Lusin theorem there is a compact set $K'\subset B$ such that
$f$ is differentiable on $K'$,
$f|_{K'}$, $L|_{K'}$ are continuous, and $|B\setminus K'|<\eps/2$. Hence
\begin{equation}
\label{Beq2}
\lim_{{\scriptstyle K'\ni y\to x}\atop {\scriptstyle y\neq x}}
\frac{|f(y)-f(x)-L(x)(y-x)|}{|y-x|} =0
\quad
\mbox{for all $x\in K'$.}
\end{equation}
This condition is however, weaker than the one required in the Whitney theorem -- we need uniform
convergence over a compact set as $|x-y|\to 0$. Let
$$
R(x,y)=\frac{|f(y)-f(x)-L(x)(y-x)|}{|y-x|}
$$
and let
$$
\eta_k(x) = \sup\{R(x,y):\, K'\ni y\neq x,\ |x-y|<1/k\}
\quad
k=1,2,\ldots
$$
Note that the $\eta_k$ are measurable. Also, condition \eqref{Beq2} means that for every $x\in K'$,
$\eta_k(x)\to 0$ as $k\to\infty$.  According to Egorov's theorem
there is another compact set $K\subset K'$ such that $|K'\setminus K|<\eps/2$
and
$$
\eta_k\rightrightarrows 0
\quad
\text{uniformly on $K$ as $k\to\infty$.}
$$
Hence
$$
\lim_{{\scriptstyle x,y\in K, x\neq y}\atop {\scriptstyle |x-y|\to 0}}
\frac{|f(y)-f(x)-L(x)(y-x)|}{|y-x|} =0.
$$
According to the Whitney Theorem~\ref{BT6}, there is $F\in C^1(\bbbr^n)$ such that $F|_K=f|_K$,
$DF|_K=Df|_K$, $|B\setminus K|<\eps$. Multiplying $F$ by a function $\varphi\in C_0^\infty(B)$ that is equal to $1$ in a neighborhood of $K$
we may further assume that $F$ has compact support in $B$. Since both functions $f$ and  $F$ and their derivatives
vanish outside $B$, we have that $F=f$ and $DF=Df$ on $\bbbr^n$ except for a set of measure less than $\eps$.

The general case follows from a partition of unity argument.
\begin{lemma}[Partition of unity]
\label{WT1-T3}
Let $\Omega\subset\bbbr^n$ be open. Then there is a family of balls $B(x_i,r_i)\subset\Omega$, $r_i\in (0,1)$,
$i=1,2,\ldots$ and a family of functions $\varphi_i\in C_0^\infty(B(x_i,2r_i))$ such that
\begin{enumerate}
\item
$\bigcup_{i=1}^\infty B(x_i,r_i)=\Omega$;
\item
$\bar{B}(x_i,2r_i)\subset\Omega$;
\item
Every compact set $K\subset\Omega$ intersects only finitely many balls $B(x_i,2r_i)$.
\item
$
\sum_{i=1}^\infty \varphi_i(x)=1
$
for every $x\in\Omega$.
\end{enumerate}
\end{lemma}
We will not prove this lemma.

We now prove the general case in the proof of Theorem \ref{BT5}. Let $A\subset \Omega$ be a closed set such that
$|\Omega\setminus A|<\eps/2$.
Let $\{\varphi_i\}_{i=1}^{\infty}$ be the partition of unity from Lemma~\ref{WT1-T3}.
Let $I\subset\bbbn$ be the set of all indices $i$ such that $B(x_i,2r_i)\cap A\neq\varnothing$.

For $i\in I$, set $f_i=\varphi_i f$. Clearly, we can regard $f_i$ as a
Lipschitz function on $\mathbb R^n$ that vanishes outside $B(x_i,2r_i)$. Applying the compactly supported case to
$f_i$, we find $F_i\in C^1_0(B(x_i,2r_i))$ and measurable sets
$Z_i\subset \mathbb R^n$ such that
$\sum_{i\in I}|Z_i|<\eps/2$, and
\[
F_i=f_i, \text{ and } DF_i=Df_i
\quad\text{on } \mathbb R^n\setminus Z_i .
\]
Define $F=\sum_{i\in I}F_i$.
Since the sum is locally finite in $\mathbb R^n$, we have that $F\in C^1(\mathbb R^n)$.
Note that for $x\in A$,
$$
\sum_{i\in I}\varphi_i(x)= \sum_{i=1}^\infty\varphi_i(x)=1
\qquad
\text{and}
\qquad
\sum_{i\in I}D\varphi_i(x)= \sum_{i=1}^\infty D\varphi_i(x)=0.
$$
This easily implies that
$$
Df(x)=DF(x) \ \ \text{and}\ \ f(x)=F(x)
\ \
\text{on } A\setminus Z, \ \text{ where } \ Z=\bigcup_{i\in I} Z_i.
$$
Since $|\Omega\setminus(A\setminus Z)|\leq |\Omega\setminus A|+|Z|<\eps$, there is a closed set $E\subset A\setminus Z$ satisfying $|\Omega\setminus E|<\eps$. Thus $f_\eps=F$ satisfies the first part of the claim of Theorem \ref{BT5}.

Finally, assume that $|f|\leq M$. Let $\psi\in C^\infty(\bbbr)$ be such that $\psi(x)=x$ if $|x|\leq \tfrac{3}{2}M$, $\psi(x)=-2M$ if $x\leq -2M$ and $\psi(x)=2M$ if $x\geq 2M$. Let $F\in C^1(\bbbr^n)$ be the function constructed above, so that $F$ and $DF$ coincide with $f$ and $Df$, respectively, on a closed $E\subset \Omega$, with $|\Omega\setminus E|<\eps$. Set $f_\eps=\psi\circ F$. Clearly, $f_\eps=f$ and $D(f_\eps)=Df$ on $E$ and $|f_\eps|\leq 2M$.
\end{proof}

\end{document}
